% Generated draft: requires mathematical and source review.
\documentclass[11pt]{article}
\usepackage[margin=25mm]{geometry}
\usepackage{fontspec}
\setmainfont{Latin Modern Roman}
\usepackage{amsmath,amssymb,mathtools,unicode-math}
\setmathfont{Latin Modern Math}
\usepackage{enumitem,needspace,xurl,hyperref}
\setlength{\parindent}{0pt}
\setlength{\parskip}{5pt}
\setlength{\emergencystretch}{3em}
\newcommand{\N}{\mathbb{N}}
\newcommand{\Z}{\mathbb{Z}}
\newcommand{\Q}{\mathbb{Q}}
\newcommand{\R}{\mathbb{R}}
\newcommand{\C}{\mathbb{C}}
\newcommand{\F}{\mathbb{F}}
\newcommand{\A}{\mathbb{A}}
\newcommand{\PP}{\mathbb{P}}
\newcommand{\E}{\mathbb{E}}
\newcommand{\eps}{\varepsilon}
\newcommand{\dd}{\,\mathrm{d}}
\begin{document}
\section*{UnsolvedMath ID: 30004259}
\textbf{Source problem code:} catalogue entry
\subsection{Definitions and mathematical statement}
The supplied record uses \(\mathrm{FO}_n\) in its cleaned statement and \(F O_n^+\) in its original statement, without defining either notation. The following formulation adopts the conventional free orthogonal quantum group interpretation suggested by the title; this interpretation requires human verification against the cited report.

For each positive integer \(n\), let \(A_n=C^u(O_n^+)\) be the universal unital complex \(C^*\)-algebra generated by elements \(u_{ij}=u_{ij}^*\), \(1\leq i,j\leq n\), satisfying
\[
\sum_{k=1}^n u_{ik}u_{jk}=\delta_{ij}1,
\qquad
\sum_{k=1}^n u_{ki}u_{kj}=\delta_{ij}1
\quad (1\leq i,j\leq n).
\]
Equip this algebra with the coproduct
\[
\Delta_n(u_{ij})=\sum_{k=1}^n u_{ik}\otimes u_{kj}.
\]
Let \(h_n\) be its Haar state, characterized by
\[
(h_n\otimes\mathrm{id})\Delta_n(a)
=(\mathrm{id}\otimes h_n)\Delta_n(a)
=h_n(a)1
\quad(a\in A_n),
\]
and let \(\pi_{h_n}\) be its GNS representation. Define the reduced algebra
\[
B_n=\overline{\pi_{h_n}(A_n)}^{\lVert\cdot\rVert}.
\]
Under the stated interpretation, \(\mathrm{FO}_n\) denotes the discrete quantum group dual to \(O_n^+\), and the source expression \(C_r^*(\mathrm{FO}_n)\) denotes \(B_n\).

\textbf{Question.} Given \(M,N\in\Z_{\geq 1}\) with \(M\neq N\), determine whether there exists a complex \(C^*\)-algebra isomorphism
\[
\Phi:B_M\longrightarrow B_N.
\]
Here an isomorphism is a bijective complex-linear map preserving multiplication and the involution; no compatibility with the coproduct is imposed. The question concerns the underlying reduced \(C^*\)-algebras. All positive ranks, including \(1\) and \(2\), are retained as supplied. The record does not clarify whether an additional rank restriction was intended or reconcile its two quantum group notations.
\subsection{Short English statement}
For distinct positive integers \(M\) and \(N\), determine whether the reduced \(C^*\)-algebras associated with the free orthogonal quantum groups of those ranks are isomorphic. This formulation uses an explicitly stated interpretation of the source notation. The supplied open label is dated curation, not an independently verified assertion that every such pair remains unresolved.
\subsection{Sources}
\begin{itemize}
\item Supplied local UnsolvedMath registry record, permanent ID \(30004259\), source problem code \texttt{catalogue entry}, in the Oberwolfach Reports Open Problems collection. No individual proposer is supplied; the proposed year is recorded as 2019.
\item Supplied report citation: \emph{Mini-Workshop: Operator Algebraic Quantum Groups (2020)}, Oberwolfach Reports, DOI \texttt{10.4171/owr/2019/45}, pp.~2821--2867. \url{https://doi.org/10.4171/owr/2019/45}. The displayed year 2020 and the 2019 provenance are retained as supplied.
\item Supplied literature source labelled ``OWR 2019 45'': \url{https://ems.press/content/serial-article-files/46823}.
\item Supplied literature source labelled ``Voigt'': \url{https://arxiv.org/abs/0911.2999}. No further bibliographic details are supplied.
\item Literature assessment date: \texttt{2026-08-22}. This is supplied dated curation, not independently verified here. The supplied status is ``open''. The assessment states that \(K\)-theory computations for reduced free-group \(C^*\)-algebras do not recover the rank parameter and that the report leaves the distinct-rank isomorphism question open. Its reference to free-group algebras is retained as a source qualification; its applicability to the free orthogonal quantum group question requires review.
\item The record labels the statement \texttt{corrected\_verified} and reports that it was checked against the official Oberwolfach report. This is a supplied verification claim, not a verification performed for this draft. The background expressly requests independent verification of time-sensitive status and cited claims.
\end{itemize}
\end{document}
